% Standalone version — no references to prior work framing
% Model: GPT-4o | Embedding: BGE-M3 | k_rules=5, n_rounds=5, n_seeds=3

\section{Methodology}
\label{sec:methodology}

This section describes the methodology of an LLM-based intrusion detection
system (IDS) that generates interpretable, threshold-based policy rules for
classifying IoT network traffic. The system is evaluated across three
dimensions: (i)~multi-class attack-family classification
(\S\ref{subsec:multiclass}), (ii)~zero-day anomaly detection via Suspicious
Multi-Class Confidence (SMCC) scoring (\S\ref{subsec:zeroday}), and
(iii)~an anonymization ablation pipeline for privacy-utility tradeoff analysis
(\S\ref{subsec:anon}). Together, these components constitute a privacy-aware,
open-set intrusion detection system.

\subsection{Multi-Class Classification}
\label{subsec:multiclass}
%-----------------------------------------------------------------------

The system retrieves semantically similar training flows from a Chroma vector
store and prompts GPT-4o to emit a class label together with a natural-language
justification rule. Multi-class classification is realised through three
components.

\textbf{Label encoding.} Each dataset is labelled with its full taxonomy of
attack families (e.g.\ DDoS, DoS, Mirai, MITM-ArpSpoofing, Recon, and
BenignTraffic for CIC-IoT2023). Labels are stored alongside feature vectors in
the vector store so that retrieved neighbours carry class provenance.

\textbf{Prompt design.} The classification prompt includes the complete set of
$K$ known attack classes, instructing the LLM to assign one of these class
labels. Majority voting across the top-$k$ retrieved neighbours' labels
provides a soft prior that is passed to the LLM as in-context evidence.

\textbf{Per-class rule head.} Following classification, the rule-generation
module synthesises one decision-rule vector per detected class. Each rule vector
consists of $k_{\text{rules}} = 5$ threshold conditions expressed over the
most discriminative features selected by permutation-importance ranking
($n_{\text{top}} = 5$). Rules are emitted in a structured JSON format
(\texttt{feature\_name}, \texttt{value}, \texttt{op}) to allow downstream
deployment as ACL-style firewall entries.

%-----------------------------------------------------------------------
\subsection{Zero-Day Detection Protocol}\label{subsec:zeroday}
%-----------------------------------------------------------------------

Rule-based IDS systems are inherently closed-set by design and cannot identify
a truly novel attack class by name. This evaluation serves a different purpose.
We assess whether the LLM's retrieval confidence degrades gracefully when
presented with out-of-distribution traffic, enabling an anomaly flag even in
the absence of a matching rule.

\textbf{Holdout split construction.} For each dataset and a designated withheld
class $c^{*}$, all flows belonging to $c^{*}$ are removed from the training
split and from the Chroma vector store. The remaining $K-1$ known classes form
the closed-set training regime. At inference time, $c^{*}$ flows are injected
into the test stream alongside balanced samples from the known classes.

\textbf{SMCC thresholding.} Rather than requiring the LLM to label unseen
traffic, we monitor the \emph{maximum softmax confidence} of the retrieved
neighbour distribution. For a query $\mathbf{x}$, let $p_j$ denote the
proportion of retrieved neighbours labelled with class $j$. The SMCC score is
defined as:
\[
  \mathrm{SMCC}(\mathbf{x}) = \max_{j \in \mathcal{K}} p_j ,
\]
where $\mathcal{K}$ is the set of known classes. When
$\mathrm{SMCC}(\mathbf{x}) < \tau$, the query is flagged as a potential
zero-day event. The threshold $\tau$ is swept from 0 to 1 and the
\emph{Zero-Day Detection Rate} (ZDR) and \emph{False Positive Rate} (FPR) are
recorded at each point, yielding the threshold-sensitivity curves
reported in \S\ref{subsec:threshold}.

\textbf{Zero-Day Detection Rate.} ZDR is defined as the fraction of withheld-class
samples correctly flagged as anomalous:
\[
  \mathrm{ZDR} = \frac{\text{TP}_{c^{*}}}{\left| \mathcal{D}_{c^{*}} \right|} ,
\]
where $\mathcal{D}_{c^{*}}$ is the set of test flows belonging to the withheld
class.

%-----------------------------------------------------------------------
\subsection{Anonymization Pipeline Design}\label{subsec:anon}
%-----------------------------------------------------------------------

A key security concern for LLM-based IDS is that the prompt necessarily exposes
raw network feature names (e.g.\ \texttt{Header\_Length}, \texttt{SrcAddr}),
which could leak deployment-sensitive topology or vendor information. The
anonymization ablation evaluates whether replacing semantically informative
feature names with opaque identifiers ($f_0, f_1, \ldots, f_N$) degrades
detection performance.

\textbf{Anonymization map.} For each dataset, a bijective mapping
$\sigma: \text{feature\_name} \to f_i$ is constructed once and held fixed
across all seeds. The map covers all features in the dataset; the LLM
receives only the anonymized identifiers and numeric values.

\textbf{Ablation design.} Two conditions are compared across $n_{\text{seeds}}=3$
random seeds and $n_{\text{rounds}}=5$ iterative rule-refinement rounds:
(a)~\emph{Original}: features presented with their semantic names; and
(b)~\emph{Anonymized}: features presented with opaque identifiers.

\textbf{Privacy-utility tradeoff metric.} Performance is reported as
macro-averaged F1 ($\overline{F1}$) on a held-out test set. The utility
cost of anonymization is quantified as $\Delta F1 = \overline{F1}_{\text{anon}}
- \overline{F1}_{\text{orig}}$. Feature-set consistency across seeds is measured
by the Jaccard index over the selected $k_{\text{rules}} = 5$ features:
\[
  J = \frac{\left| \mathcal{F}_{\text{orig}} \cap \mathcal{F}_{\text{anon}} \right|}
           {\left| \mathcal{F}_{\text{orig}} \cup \mathcal{F}_{\text{anon}} \right|} .
\]

%-----------------------------------------------------------------------
\subsection{Robustness Analysis}\label{subsec:smcc}
%-----------------------------------------------------------------------

To evaluate robustness of the generated rules against an adversary who is aware
of the rule structure (Kerckhoffs' principle), we conduct an
\emph{Exploitability Sweep Rate} (ESR) experiment. For a given perturbation
budget $\varepsilon$ (expressed as a fraction of each feature's interquartile
range), an adversary perturbs withheld-class flows toward the nearest
known-class centroid. We measure the fraction of originally-flagged flows that
evade detection after perturbation:
\[
  \mathrm{ESR}(\varepsilon) = \frac{\text{flows that evade at budget } \varepsilon}
                                   {\text{flows flagged at } \varepsilon = 0} .
\]
This analysis is performed for five detection conditions: LLM with original
feature names (\texttt{LLM\_orig}), LLM with anonymized names
(\texttt{LLM\_anon}), Decision Tree depth-3 (\texttt{DT\_d3}), Decision Tree
depth-5 (\texttt{DT\_d5}), and Random Forest depth-5 (\texttt{RF\_d5}).
The critical budget $\varepsilon_{0.5}$ at which $\mathrm{ESR} = 0.5$ quantifies
the adversarial hardness of each detector. Experiments are repeated for three
seeds (42, 123, 456) and the mean and standard deviation of $\varepsilon_{0.5}$
are reported in Table~\ref{tab:robustness}.
